\section{Introduction}
\label{sec:introduction}

Coronary artery disease (CAD) remains the leading cause of mortality worldwide, accounting for approximately 9 million deaths annually~\cite{who2023cvd}. Early and accurate identification of patients at risk is central to clinical management, since timely intervention with lifestyle modification, pharmacotherapy, or revascularisation can substantially alter long-term outcomes. The conventional diagnostic pathway combines clinical history, non-invasive testing (resting and exercise ECG, thallium imaging), and ultimately invasive coronary angiography. While angiography remains the gold standard, its cost, invasiveness, and limited availability motivate the development of decision-support tools that can identify which patients are most likely to benefit from referral.

Machine learning has been applied extensively to this task, with the UCI Heart Disease dataset~\cite{detrano1989cleveland} serving as one of the most widely studied benchmarks in medical machine learning. Despite a sample size of only a few hundred patients, the dataset captures the heterogeneity of routine cardiology workup --- demographic, symptomatic, biochemical, and imaging-derived features --- and remains relevant for evaluating modern modelling pipelines under realistic small-data constraints. Performance estimates reported on this dataset have ranged widely~\cite{detrano1989cleveland,gennari1989models}, reflecting differences in preprocessing, validation strategy, and the level of methodological rigour applied to avoid optimistic bias.

The central methodological challenge in clinical machine learning on small datasets is producing a generalisation estimate that is both unbiased and reproducible. Standard $k$-fold cross-validation is well known to be biased upward when used simultaneously for hyperparameter selection and performance reporting, because the same data informs both decisions~\cite{cawley2010overfitting}. Repeated nested cross-validation (rnCV)~\cite{varma2006bias,krstajic2014cross} addresses this by introducing a strict separation between an outer loop for performance estimation and an inner loop for hyperparameter tuning, with the procedure repeated multiple times to reduce variance. Combined with strict containment of all preprocessing and feature-selection steps inside the cross-validation folds, rnCV provides a defensible workflow for reporting generalisation performance on small datasets.

This work presents a complete rnCV-based machine learning pipeline for binary heart disease classification on the Cleveland subset of the UCI dataset. The pipeline is implemented following object-oriented design and addresses the assignment requirements through five specific contributions:

\begin{enumerate}
    \item \textbf{Exploratory data analysis.} A systematic EDA characterises the dataset's structure, identifies the two missing values, the cholesterol outlier, the mild class imbalance, and the most class-discriminative features. Findings inform every subsequent preprocessing decision.

    \item \textbf{Object-oriented rnCV implementation.} A reusable \texttt{RepeatedNestedCV} class encapsulates the full pipeline, accepting a list of estimators with their hyperparameter spaces. Stratified folds, per-repetition seeding, and strict in-fold preprocessing prevent leakage; non-parametric bootstrap CIs quantify uncertainty around median metrics.

    \item \textbf{Algorithm comparison and winner selection.} Seven algorithms spanning linear, probabilistic, discriminant, and ensemble paradigms are compared under both default and tuned hyperparameters. Performance is reported on a battery of clinically aligned metrics including MCC, ROC-AUC, PR-AUC, balanced accuracy, F1, recall, specificity, and precision. Linear Discriminant Analysis is identified as the winner.

    \item \textbf{Stable feature selection.} Permutation-importance feature selection inside the outer loop identifies a five-feature stable subset (\texttt{thal}, \texttt{ca}, \texttt{cp}, \texttt{sex}, \texttt{exang}) that matches or improves the full-feature performance while substantially reducing the model's input requirements --- a result of practical importance for deployment.

    \item \textbf{Final model and SHAP interpretation.} A complete deployable pipeline is trained on all available data and serialised. Post-hoc SHAP analysis~\cite{lundberg2017unified} confirms that the model's decision logic aligns with established cardiology knowledge, including the well-documented silent-ischaemia signature in the chest-pain feature.
\end{enumerate}

A bonus error analysis on a held-out validation split additionally identifies the dominant error mechanism (over-reliance on the thallium scan when it is atypical for the patient's true disease status) and produces a clinician-oriented summary suitable for guiding deployment decisions.

The remainder of this report is structured as follows. Section~\ref{sec:methods} describes the dataset, preprocessing, rnCV pipeline, and evaluation methodology. Section~\ref{sec:results} reports EDA findings, algorithm comparison, feature selection, final model performance with SHAP interpretation, and the bonus error analysis. Section~\ref{sec:conclusions} summarises the main findings, discusses limitations, and outlines directions for further investigation.
